\section{Conclusion}
\label{sec:conclusion}

This survey has traced the evolution of Novel View Synthesis from its origins
in image-based rendering through the neural radiance field era, the Gaussian
Splatting revolution, the emergence of generalisable feed-forward reconstruction,
and the convergence with generative world modeling.
The organizing perspective we have pursued---that \nvs{} provides essential geometric building blocks for the visual-generative branch of
inference engine of Generative World Modeling---finds strong empirical support
in the recent literature, where the most capable world models (Genie~2,
Cosmos, ViewCrafter) are distinguished from their predecessors precisely by
their explicit incorporation of spatial consistency constraints derived from
the NVS research tradition.

The unified continuous-spectrum taxonomy we have proposed organises more than
180 methods along three axes (input regime, temporal scope, controllability)
and makes explicit their world-modeling implications.
Viewed through this lens, the history of the field is a systematic expansion
of coverage along the spectrum: from dense multi-view static reconstruction
at the origin, through sparse-view and single-image generalisation in the
middle range, toward the 4-D, action-conditioned, physically plausible
generation that defines the far end of the spectrum and the open research frontier.

The field is at a genuine inflection point.
The spatial grounding developed by the NVS community and the temporal and
action capabilities developed by the world modeling community are actively
converging, as evidenced by systems such as Genie~2, Cosmos, ViewCrafter,
and the neural driving simulators reviewed in Section~\ref{sec:world}.
This convergence is not merely a metaphor: architecturally, the most capable
recent systems are genuine fusions of explicit 3-D representations and
large-scale video generation priors, combining the geometric accuracy of
the former with the coverage and controllability of the latter.

Closing the geometry--generation trade-off, achieving long-horizon temporal
coherence, and grounding predictions in physical laws are the three primary
technical challenges ahead.
The five research directions proposed in Section~\ref{sec:open} represent
our assessment of the most productive paths to addressing these challenges.
Progress along these directions will not only advance the field of \nvs{}
and world modeling in isolation, but will also enable downstream capabilities
in autonomous driving, robotic manipulation, and interactive AI agents that
require the ability to reason about and generate plausible futures of the
physical world.

The survey is an ongoing work.
Content will be further expanded and updated as the field continues to
evolve at its current rapid pace.

% ─────────────────────────────────────────────────────────────────────────────

% References are managed in references.bib (see main.tex)
\bibliographystyle{plainnat}
\bibliography{references}


\end{document}
